\scanpage{28}
\textbf{C. R. Acad. Sc. Paris, t. 292 (19 janvier 1981)\hfill Série I --- }\mspage{213}

TOPOLOGIE. --- \textit{Théorie axiomatique de l'orientation des variétés topologiques à bord.} Note (*) de \textbf{Marcus Wanderley}, présentée par Henri Cartan,

On définit ce qu'on entend par une théorie de l'orientation des variétés topologiques (à bord) de dimension $\leqq\mathrm{N}$, et on prouve l'unicité (à isomorphisme près) d'une telle théorie lorsque $\mathrm{N}\ne4$, en s'appuyant sur un résultat de Kirby (pour $n\ne4$, tout homéomorphisme de $\mathbf{R}^n$ de degré $+1$ est « stable »).

\textit{One defines what is understood by an orientation theory for the topological manifolds (with boundary) of dimension $\leqq\mathrm{N}$. One proves that such a theory is unique (up to isomorphism) if $\mathrm{N}\ne4$, using a result of Kirby (for $n\ne4$ every homeomorphism of $\mathbf{R}^n$ of degree $+1$ is ``stable'').}

\textsc{Notations}. --- Dans cette Note, $\mathrm{N}$ désigne un entier positif ou le symbole $\infty$. On désigne par $\mathscr{V}^{\bullet}_{\leqq\mathrm{N}}$ la catégorie dont les objets sont les variétés topologiques à bord de dimension $\leqq\mathrm{N}$, les morphismes étant les plongements ouverts. La sous-catégorie pleine de $\mathscr{V}^{\bullet}_{\leqq\mathrm{N}}$ formée des variétés topologiques (i. e. à bord vide) de dimension $\leqq\mathrm{N}$ est notée par $\mathscr{V}_{\leqq\mathrm{N}}$. Si $m$ est un entier positif $\leqq\mathrm{N}$, on désigne par $\mathscr{V}^{\bullet}_m$ (resp. $\mathscr{V}_m$) la sous-catégorie pleine de $\mathscr{V}^{\bullet}_{\leqq\mathrm{N}}$ (resp. $\mathscr{V}_{\leqq\mathrm{N}}$) des variétés pures de dimension $m$. Lorsque l'on a un foncteur contravariant :
\[
\mathrm{F}:\quad\mathscr{V}^{\unsure{0|\mathrm{o}}}_{\leqq\mathrm{N}}\to\mathrm{Ens},
\]
on le prolonge à un foncteur :
\[
\mathrm{F}^{\bullet}:\quad\mathscr{V}^{\bullet\;0}_{\leqq\mathrm{N}}\to\mathrm{Ens},
\]
par la formule :
\[
\mathrm{F}^{\bullet}(\mathrm{V})=\mathrm{F}(\operatorname{Int}(\mathrm{V}))\qquad [\mathrm{V}\in\operatorname{ob}(\mathscr{V}^{\bullet}_{\leqq\mathrm{N}})].
\]

\textsc{Définition}. --- Une théorie $\mathrm{T}=(\Omega,\delta,\nu)$ de l'orientation des variétés topologiques de dimension $\leqq\mathrm{N}$ consiste en les données (1), (2) et (3) ci-dessous, astreintes aux axiomes $(a)$, $(b)$, $(c)$ et $(d)$ ci-dessous.

\textsc{Données}. --- (1) (\textit{Foncteur}). On se donne un foncteur contravariant :
\[
\Omega:\quad\mathscr{V}^{0}_{\leqq\mathrm{N}}\to\mathrm{Ens},
\]
appelé \textit{foncteur orientation}. Le prolongement $\Omega^{\bullet}$ de $\Omega$ à $\mathscr{V}^{\bullet\;0}_{\leqq\mathrm{N}}$ est appelé également \textit{foncteur orientation}. Si $\mathrm{V}$ est un objet de $\mathscr{V}^{\bullet}_{\leqq\mathrm{N}}$ alors $\Omega^{\bullet}(\mathrm{V})$ est appelé \textit{l'ensemble des orientations de} $\mathrm{V}$ (au sens de $\mathrm{T}$).

(2) (\textit{Stokes}). On se donne pour tout objet $\mathrm{V}$ de $\mathscr{V}^{\bullet}_{\leqq\mathrm{N}}$ une application :
\[
\delta_{\mathrm{V}}:\quad\Omega^{\bullet}(\mathrm{V})\to\Omega\,(\partial\mathrm{V}),
\]
fonctorielle en $\mathrm{V}$ pour $\mathrm{V}$ variable, appelée \textit{application de Stokes}. Si $\omega$ est une orientation de $\mathrm{V}$ alors $\delta_{\mathrm{V}}(\omega)$ est appelée \textit{orientation unduite sur le bord} $\partial\mathrm{V}$ \textit{de} $\mathrm{V}$.

(3) (\textit{Normalisation}). On se donne pour la variété $\mathrm{P}$ réduite à un point une bijection :
\[
\nu:\quad\{\pm1\}\overset{\sim}{\to}\Omega(\mathrm{P}),
\]
appelée \textit{application de normalisation}.

\textsc{Axiomes}. --- $(a)$ (\textit{Exactitude}). $\Omega$ est un « faisceau », c'est-à-dire pour tout $\mathrm{V}$ objet de $\mathscr{V}_{\leqq\mathrm{N}}$ la \textit{restriction} $\Omega_{\mathrm{V}}$ de $\Omega$ à la catégorie des ouverts de $\mathrm{V}$ avec comme morphismes les inclusions, est un faisceau (\textit{cf.} [1]).

% archive sheet number: 27
